\section{Introdução}

Requisitos escritos em linguagem natural (LN) irrestrita permanecem o veículo predominante para comunicação entre partes interessadas em projetos industriais de software e sistemas embarcados~\cite{darif2026controlled}. Entretanto, a flexibilidade inerente à prosa livre frequentemente introduz ambiguidades lexicais, omissão de atores, voz passiva e agrupamento de múltiplas obrigações independentes em uma única sentença, prejudicando a rastreabilidade, a verificação formal e a garantia de qualidade em escala~\cite{incose2023guide, grosser2024benchmarking}. Para mitigar esses problemas sem exigir linguagens lógicas matemáticas inacessíveis a clientes não técnicos, sistemas de templates semiformais --- como EARS e MASTeR (\textit{Muster für Anforderungen und System-Tests im Engineering} / SOPHIST)~\cite{sophist2020master, rupp2014anforderungsschablonen} --- estabelecem moldes sintáticos fixos que elevam a clareza, a singularidade e a conformidade normativa~\cite{grosser2023comparative, grosser2024benchmarking}.

Apesar de seus benefícios comprovados em domínios críticos de segurança (como os setores automotivo e aeroespacial), migrar especificações legadas contendo milhares de requisitos em texto livre para o padrão MASTeR impõe alto custo manual: exige domínio simultâneo do negócio e da gramática de sub-templates, além da decomposição (\textit{splitting}) frequente de requisitos compostos em múltiplos requisitos atômicos~\cite{rosskothen2026converting}. Recentemente, Ro\ss kothen et al.~\cite{rosskothen2026converting} apresentaram uma arquitetura orientada à indústria para automatizar essa conversão em lote utilizando Grandes Modelos de Linguagem (LLMs) organizados em um ciclo reflexivo (\textit{Reflection Loop}~\cite{ji2023mitigating}), combinando um validador sintático MASTeR, um avaliador alinhado ao guia INCOSE~\cite{incose2023guide} e um filtro de acionabilidade de feedback.

Uma análise detalhada da arquitetura e do pacote de replicação público de Ro\ss kothen et al.~\cite{rosskothen2026zenodo}, contudo, expõe três limitações críticas para adoção em larga escala:
\begin{enumerate}
    \item \textbf{Sobrecarga multiplicativa de chamadas generativas:} Para cada requisito fonte decomposto em $k$ requisitos candidatos, uma única iteração de validação executa até $2k + 1$ chamadas autorregressivas a LLMs juízes (\textit{Template Validator}, \textit{INCOSE Validator} em nove categorias e \textit{Feedback Filter}), reenviando repetidamente prompts extensos (~1.400 tokens) apenas para produzir justificativas textuais intermediárias;
    \item \textbf{Paradoxo de \textit{over-editing} sem calibração:} Nos dados empíricos publicados pelos próprios autores~\cite{rosskothen2026zenodo}, o modelo \texttt{Qwen3-235B} sem validadores generativos (\texttt{minimal-qwen235}) obteve 68\% de avaliações humanas positivas (categorias C1/C2) e pontuações automáticas ligeiramente superiores às da configuração completa com juízes (\texttt{full-qwen235}: \texttt{CorpBERTF1} de $0{,}834$ contra $0{,}823$; \textit{Compound Score} de $0{,}811$ contra $0{,}797$), indicando que ciclos reflexivos não calibrados podem degradar rascunhos inicialmente corretos;
    \item \textbf{Falta de roteamento explícito por incerteza:} Embora os autores demonstrem que tanto os LLMs quanto as métricas automáticas sofrem forte queda de desempenho em requisitos complexos (\textit{split factor} $\ge 2$ ou mais de 30 palavras), o fluxo original não fornece probabilidades calibradas para decidir automaticamente quando aceitar um requisito, quando reescrevê-lo ou quando escalá-lo para inspeção humana.
\end{enumerate}

Neste artigo, investigamos a substituição da pilha de validadores generativos (\textit{LLM-as-a-Judge}) por um modelo discriminativo de decisão rápida (\textit{System One}), especificamente o \textbf{JEV}~\cite{typesafe2026jev}. Diferentemente de LLMs autorregressivos que decodificam tokens de texto sequencialmente para justificar notas, o JEV avalia um estado estruturado contra múltiplas consultas tipadas (\texttt{Choice}, \texttt{Noul} e \texttt{Score}) em um único passe paralelo, retornando distribuições de probabilidade calibradas em milissegundos. Combinando validação discriminativa no ciclo reflexivo com o padrão arquitetural \textit{Accept-Rewrite-Escalate}, propomos o \textbf{JEV-Cascade Reflection Loop}.

As principais contribuições deste trabalho são:
\begin{itemize}
    \item \textbf{(i)} O projeto arquitetural do \textit{JEV-Cascade Reflection Loop}, que condiciona a validação simultaneamente ao requisito fonte $S$ e ao candidato $R_i$, eliminando por construção o estágio extra de filtragem de feedback por LLM e mapeando falhas para instruções determinísticas de reparo;
    \item \textbf{(ii)} Um mecanismo de roteamento triplo baseado em probabilidades calibradas que protege rascunhos de alta confiança contra \textit{over-editing} (\textit{Fast Accept}), aciona reescrita apenas diante de violações claras (\textit{Guided Rewrite}) e isola casos de incerteza epistêmica para engenheiros de requisitos (\textit{Escalate});
    \item \textbf{(iii)} Um protocolo e repositório experimental reprodutível avaliado sobre os 249 requisitos fonte e 410 requisitos MASTeR de referência~\cite{grosser2023comparative, rosskothen2026zenodo}, contrastando eficácia sintático-semântica, latência e custo operacional em modelos abertos (\texttt{Qwen3-235B}) e comerciais (\texttt{Gemini 3.5 Flash}).
\end{itemize}
